\documentclass{article}
\usepackage{graphicx} % Required for inserting images

\title{Ultrahangos mérés}
\author{Kern Luca, Csongor Vincze}
\date{February 2026}

\begin{document}

\maketitle

\section{1. Feladat Ismerkedés az oszcilloszkóp és a függvénygenerátor használatával.}

Beállítottuk a 4 V amplitúdójú, 1 V eltolású 115 Hz-es szinuszjelet. Csatlakoztattuk a jelet az oszcilloszkóphoz is. Beállítottuk a triggert, és stabilizáltuk, ízlésesen a képernyő közepér hoztuk a jelet.
Végül közelítőleg lemértük a jel amplitúdóját és frekvenciáját.

\begin{table}[h!]
    \centering
    \begin{tabular}{|l|c|c|}
        \hline
        Paraméter & Beállított érték & Mért érték \\ \hline
        Amplitúdó & ?? & ?? \\ \hline
        Frekvencia & ?? & ?? \\ \hline
    \end{tabular}
    \caption{A beállított és a mért amplitúdó és frekvencia értékek.}
    \label{tab:amplitudo_frekvencia}
\end{table}

A különbség az RMS és a csúcstól csúcsig mért értékek között a fő különbség, hogy a csúcstól csúcsig mért érték az amplitúdó kétszerese, míg az effektív érték a jel négyzetének a kiátlagolásának a gyöke.

Megváltoztattuk a csatolási módot DC-ről AC-re, így kiküszöböltük a konstans tagú feszültség miatti eltolást.
Ha a csatolást "Ground"-ra állítottuk a következő történt:

Végül a kurzor egyenesek segítségével is megmértük az amplitúdó és frekvencia értékeit a jelnek.


\begin{figure}[htp]
    \centering
    \includegraphics[width=4cm]{your_image_here.png}
    \caption{Az oszcilloszkóp különböző funkcióinak kipróbálása}
    \label{probalkozas}
\end{figure}

\section{Ultrahangos adó-vevő páros átviteli karakterisztikája
a frekvencia függvényében}

Lehelyeztük a két adó-vevőt ??? cm távolságra, és kissé elforgattuk őket, hogy ne zavarjanak be a visszaverődő hullámok. Ezután beállítottuk az 5,9 V-os eltolásmentes színuszjelet. Csatlakoztattuk az koax kábelt a jelgenerátorhoz, aztán egy koaxkábeles elosztóval bekötöttük a jelet az oszcilloszkópba és az adóba is. Az oszcilloszkóp másik bemenetéhez pedig csatlakoztattuk a vevőt. Ezután beállítottuk a triggert.

\begin{table}[h!]
    \centering
    \begin{tabular}{|c|c|}
        \hline
        Frekvencia [kHz] & Amplitúdó [V] \\ \hline
        1 & \\ \hline
        5 & \\ \hline
        10 & \\ \hline
        15 & \\ \hline
        20 & \\ \hline
        25 & \\ \hline
        30 & \\ \hline
        35 & \\ \hline
        40 & \\ \hline
        45 & \\ \hline
        50 & \\ \hline
        55 & \\ \hline
        60 & \\ \hline
        65 & \\ \hline
        70 & \\ \hline
        75 & \\ \hline
        80 & \\ \hline
        85 & \\ \hline
        90 & \\ \hline
        95 & \\ \hline
        100 & \\ \hline
    \end{tabular}
    \caption{Az adó-vevő páros átviteli karakterisztikájának mérési adatai.}
    \label{tab:atviteli_karakterisztika}
\end{table}


\begin{enumerate}
    \item A legalacsonyabb frekvenciájú jel, amit ezzel a beállítással mérni lehet: 
    \item A frekvencia, amitől kezdve hallani kezdjük a jelet: 
    \item A legmagasabb frekvencia, amit még hallunk: 
    \item A rezonanciafrekvenciák: 
    \item A rezonanciafrekvencia, amely a legnagyobb csúcsot adja: 
\end{enumerate}


\begin{table}[h!]
    \centering
    \begin{tabular}{|c|c|}
        \hline
        Frekvencia [kHz] & Amplitúdó [V] \\ \hline
         &  \\ \hline
         &  \\ \hline
         &  \\ \hline
         &  \\ \hline
         &  \\ \hline
    \end{tabular}
    \caption{Rezonanciafrekvencia körüli adatpontok.}
    \label{rezonancia freki}
\end{table}


A rezonanciafrekvencia értéke:


\begin{table}[h!]
    \centering
    \begin{tabular}{|c|c|}
        \hline
        Frekvencia [kHz] & Amplitúdó [V] \\ \hline
         &  \\ \hline
         &  \\ \hline
         &  \\ \hline
         &  \\ \hline
         &  \\ \hline

         &  \\ \hline
         &  \\ \hline
         &  \\ \hline
         &  \\ \hline
         &  \\ \hline
    \end{tabular}
    \caption{Csúcstól csúcsig amplitúdó értékek}
    \label{Vpp}
\end{table}


\section{Hangsebesség meghatározása az átvitt jel fázisának mérésével}


\begin{table}[h!]
    \centering
    \begin{tabular}{|c|c|c|}
        \hline
        Távolság [mm] & Idő & Sebesség \\ \hline
        0 & & \\ \hline
        2 & & \\ \hline
        4 & & \\ \hline
        6 & & \\ \hline
        8 & & \\ \hline
        10 & & \\ \hline
        12 & & \\ \hline
        14 & & \\ \hline
        16 & & \\ \hline
        18 & & \\ \hline
        20 & & \\ \hline
        22 & & \\ \hline
        24 & & \\ \hline
    \end{tabular}
    \caption{A hangsebesség mérése a távolság függvényében.}
    \label{tab:hangsebesseg_meres}
\end{table}

A hullámhossz:

\begin{figure}[htp]
    \centering
    % \includegraphics[width=4cm]{your_image_here.png}
    \caption{Az oszcilloszkóp különböző funkcióinak kipróbálása}
    \label{hangsebesség_grafikon}
\end{figure}


\section{Állóhullám mintázat}



\begin{table}[h!]
    \centering
    \begin{tabular}{|c|c|c|}
        \hline
        Távolság [mm] & Amplitúdó \\ \hline
        0 &  \\ \hline
        2 &  \\ \hline
        4 &  \\ \hline
        6 &  \\ \hline
        8 &  \\ \hline
        10 &  \\ \hline
        12 &  \\ \hline
        14 &  \\ \hline
        16 &  \\ \hline
        18 &  \\ \hline
        20 &  \\ \hline
        22 &  \\ \hline
        24 &  \\ \hline
    \end{tabular}
    \caption{állóhullám amplitúdófüggese a távoldágtól}
    \label{állóhullám}
\end{table}







\end{document}
